\documentclass[a4paper, serif, 10pt]{moderncv}

%% moderncv
% color and style
\moderncvcolor{black}
\moderncvstyle{banking}

%% page layout/margins
\usepackage[top=2.5cm, bottom=2.5cm, left=1.5cm, right=1.5cm]{geometry}

\nopagenumbers{}

%% Babel config for Italian language
% ref:
% - https://latex3.github.io/babel/guides/locale-italian.html
\usepackage[italian]{babel}

%% fontspec
\usepackage{fontspec}

\IfFontExistsTF{Linux Libertine}{
  \setmainfont[Ligatures={TeX, Common}, Numbers=OldStyle]{Linux Libertine}
}{}
\IfFontExistsTF{Linux Libertine O}{
  \setmainfont[Ligatures={TeX, Common}, Numbers=OldStyle]{Linux Libertine O}
}{}

%% CTeX
% CJK support, used to show CJK characters in the resume
%
% - heading=false: disable ctex heading
% - punct=kaiming: use kaiming punctuations styles for CJK
% - scheme=plain: use plain scheme, do not override \normalsize font size
% - space=auto: space settings for CJK characters
%
% CJK fonts are resolved per locale language, e.g. Noto Serif/Sans CJK JP
% for Japanese resumes, with fallback to the next region if not installed
%
% ref:
% - http://ctan.mirrorcatalogs.com/language/chinese/ctex/ctex.pdf
\usepackage[UTF8, heading=false, punct=kaiming, scheme=plain, space=auto]{ctex}

\IfFontExistsTF{Noto Serif CJK SC}{
  \setCJKmainfont{Noto Serif CJK SC}
}{}
\IfFontExistsTF{Noto Sans CJK SC}{
  \setCJKsansfont{Noto Sans CJK SC}
}{}

%% line spacing
\usepackage{setspace}
\setstretch{1.125}

%% URL styling - use normal text instead of monospace
\urlstyle{same}

%% Patch \httplink and \httpslink to handle full URLs
% The original moderncv macros always prepend http:// or https:// to URLs.
% This causes issues when the URL already has a protocol (e.g., https://example.com)
% resulting in malformed URLs like https://https://example.com.
% This patch checks if the URL already contains "://" and uses it directly if so.
% Note: We use \str_set:Nx (x-expansion) to expand macros like \@homepage first.
\ExplSyntaxOn
\str_new:N \g_yamlresume_url_str

\RenewDocumentCommand{\httpslink}{O{}m}{
  \str_set:Nx \g_yamlresume_url_str {#2}
  \str_if_in:NnTF \g_yamlresume_url_str {://}
    { \url{#2} }
    { \url{https://#2} }
}

\RenewDocumentCommand{\httplink}{O{}m}{
  \str_set:Nx \g_yamlresume_url_str {#2}
  \str_if_in:NnTF \g_yamlresume_url_str {://}
    { \url{#2} }
    { \url{http://#2} }
}
\ExplSyntaxOff

\name{Lorenzo Bianchi}{}
\title{Ingegnere della cybersicurezza}
\phone[mobile]{+39 02 8765 4321}
\email{lorenzo.bianchi@lorenzobianchi.it}
\homepage{https://www.lorenzobianchi.it}

\address{Via Torino 45, Milano, Lombardia, Italia, 20123}{}{}

\extrainfo{{\small \faLinkedin}\ \href{https://www.linkedin.com/in/lorenzo-bianchi-cyber}{@lorenzo-bianchi-cyber} {} {} {} • {} {} {} 
{\small \faGithub}\ \href{https://github.com/lorenzobianchi}{@lorenzobianchi}}

\begin{document}

\maketitle

\section{Informazioni personali}

\cvline{}{Professionista della sicurezza informatica con oltre otto anni di esperienza nella protezione di infrastrutture critiche e dati sensibili. Specializzato in monitoraggio, risposta agli incidenti e architetture Zero Trust, unisco competenze tecniche e visione del rischio per sostenere gli obiettivi aziendali.
%
\begin{itemize}
\item Guido progetti di adeguamento a NIS2, GDPR e ISO 27001.
\item Coordino team di analisti e collaboro con sviluppo, reparti operativi e ufficio legale.
\item Progetto controlli preventivi e promuovo una cultura della sicurezza misurabile.
\end{itemize}}
\section{Formazione}

\cventry{set 2014–lug 2016}
        {Laurea magistrale, Ingegneria Informatica, Punteggio: 110/110 con lode}
        {Politecnico di Milano}
        {\url{https://www.polimi.it}}
        {}
        {\begin{itemize}
\item Tesi magistrale sulla rilevazione di intrusioni in ambienti cloud tramite tecniche di apprendimento automatico.
\item Approfondimento su crittografia applicata, sicurezza delle reti e gestione del rischio.
\item Partecipazione a laboratori di hacking etico e progetti con aziende del territorio.
\end{itemize}
\textbf{Corsi}: Sicurezza dei sistemi informativi, Crittografia applicata, Architetture di rete e protocolli, Sistemi operativi e virtualizzazione, Gestione dei rischi informatici, Diritto delle nuove tecnologie}

\cventry{set 2011–lug 2014}
        {Laurea triennale, Ingegneria Informatica, Punteggio: 108/110}
        {Università degli Studi di Bologna}
        {\url{https://www.unibo.it}}
        {}
        {\begin{itemize}
\item Basi solide di programmazione, algoritmi, reti e basi di dati.
\item Progetto finale su un sistema di autenticazione a due fattori per applicazioni web.
\item Esperienza in gruppi di lavoro e presentazioni tecniche.
\end{itemize}
\textbf{Corsi}: Fondamenti di programmazione, Algoritmi e strutture dati, Reti di calcolatori, Basi di dati, Sistemi operativi}
\section{Esperienza professionale}

\cventry{mar 2020–Presente}
        {Ingegnere della cybersicurezza senior}
        {Finanza Sicura S.p.A.}
        {\url{https://www.finanzasicura.it}}
        {}
        {\begin{itemize}
\item Guido un gruppo di sei analisti nella gestione di incidenti di sicurezza su scala nazionale.
\item Ho progettato e implementato un modello di difesa Zero Trust per 4.000 dipendenti e 120 applicazioni.
\item Ho ridotto del 35\% il tempo medio di rilevamento e contenimento delle minacce tramite regole SIEM e automazione SOAR.
\item Coordino verifiche ISO 27001 e attività di adeguamento alla direttiva NIS2.
\item Collaboro con i team di sviluppo per integrare controlli di sicurezza nella catena CI/CD.
\end{itemize}
\textbf{Parole chiave}: Gestione del rischio, SIEM, Zero Trust, Conformità, Risposta agli incidenti}

\cventry{set 2016–feb 2020}
        {Cybersecurity Analyst}
        {CloudNova S.r.l.}
        {\url{https://www.cloudnova.it}}
        {}
        {\begin{itemize}
\item Monitoraggio continuo di eventi di sicurezza con SIEM, IDS/IPS e piattaforme EDR.
\item Analisi di malware, phishing e tentativi di accesso anomalo, con redazione di rapporti tecnici.
\item Gestione delle vulnerabilità e coordinamento delle correzioni con i gruppi infrastrutturali.
\item Supporto nella definizione delle politiche di accesso e nella formazione del personale.
\end{itemize}
\textbf{Parole chiave}: Monitoraggio, Analisi delle minacce, Firewall, VPN, Vulnerabilità}

\cventry{ott 2014–ago 2016}
        {Specialista di rete junior}
        {TecnoReti S.p.A.}
        {\url{https://www.tecnoreti.it}}
        {}
        {\begin{itemize}
\item Configurazione e manutenzione di apparati di rete, firewall e VPN tra sedi.
\item Diagnosi e risoluzione di problemi di connettività e sicurezza per sedi distribuite.
\item Collaborazione alla migrazione verso un'infrastruttura di rete segmentata.
\end{itemize}
\textbf{Parole chiave}: Networking, Troubleshooting, Cisco, Segmentazione}
\section{Lingue}

\cvline{Italiano}{Madrelingua o bilingue \hfill \textbf{Parole chiave}: Madrelingua}
\cvline{Inglese}{Competenza professionale completa \hfill \textbf{Parole chiave}: IELTS 7.5, Livello C1}
\cvline{Francese}{Competenza professionale limitata \hfill \textbf{Parole chiave}: DELF B1}
\section{Competenze}

\cvline{Sicurezza informatica}{Esperto \hfill \textbf{Parole chiave}: SIEM, IDS/IPS, Penetration testing, Threat hunting, Risposta agli incidenti}
\cvline{Sicurezza cloud}{Avanzato \hfill \textbf{Parole chiave}: AWS, Azure, Kubernetes, IAM, Terraform}
\cvline{Governance e conformità}{Avanzato \hfill \textbf{Parole chiave}: ISO 27001, NIS2, GDPR, Valutazione del rischio, Verifiche di conformità}
\cvline{Programmazione e automazione}{Avanzato \hfill \textbf{Parole chiave}: Python, Bash, PowerShell, SQL}
\section{Premi}

\cventry{nov 2022}
        {Associazione Italiana per la Sicurezza Informatica}
        {Premio Innovazione in Cybersecurity}
        {}
        {}
        {Riconoscimento per il progetto di una piattaforma di condivisione di indicatori di compromissione tra piccole e medie imprese italiane, con impatto misurabile sulla riduzione degli incidenti.}
\section{Certificazioni}

\cventry{giu 2021}
        {ISC2}
        {Certified Information Systems Security Professional (CISSP)}
        {\url{https://www.isc2.org/certifications/cissp}}
        {}
        {}

\cventry{ott 2022}
        {Amazon Web Services}
        {AWS Certified Security - Specialty}
        {\url{https://aws.amazon.com/certification/certified-security-specialty/}}
        {}
        {}

\cventry{apr 2023}
        {TÜV Italia}
        {ISO/IEC 27001 Lead Auditor}
        {\url{https://www.tuv.com/italy/it/}}
        {}
        {}
\section{Pubblicazioni}

\cventry{set 2023}
        {Difendere la supply chain software: modelli e pratiche}
        {Rivista Italiana di Cybersecurity}
        {\url{https://www.rivistaitalianacybersecurity.it/articoli/supply-chain}}
        {}
        {\begin{itemize}
\item Analisi dei rischi legati a dipendenze software, fornitori e componenti open source.
\item Proposta di un modello di valutazione e monitoraggio continuo della catena di fornitura.
\item Casi di studio tratti da incidenti reali e indicazioni operative per le PMI.
\end{itemize}}
\section{Progetti}

\cventry{gen 2021–dic 2021}
        {Progetto open source per la condivisione di indicatori di compromissione tra piccole e medie imprese italiane.}
        {Piattaforma condivisa di threat intelligence}
        {\url{https://github.com/lorenzobianchi/threat-intel-condivisa}}
        {}
        {\begin{itemize}
\item Raccolta e normalizzazione di indicatori di compromissione da fonti eterogenee.
\item API per l'integrazione con SIEM e piattaforme di ticketing.
\item Documentazione e laboratori per l'adozione da parte di organizzazioni con risorse limitate.
\end{itemize}
\textbf{Parole chiave}: Threat intelligence, STIX/TAXII, Open source}

\cventry{feb 2019–lug 2019}
        {Ambiente virtualizzato per esercitazioni su vulnerabilità web e di rete.}
        {Laboratorio di penetration testing}
        {\url{https://github.com/lorenzobianchi/pentest-lab}}
        {}
        {\begin{itemize}
\item Configurazione di macchine vulnerabili e scenari di attacco controllati.
\item Sperimentazione di tecniche di enumerazione, sfruttamento e post-sfruttamento.
\item Redazione di report con raccomandazioni di correzione.
\end{itemize}
\textbf{Parole chiave}: Penetration testing, OWASP, Virtualizzazione}
\section{Interessi}

\cvline{Sicurezza informatica}{CTF, Ingegneria inversa, OSINT}
\cvline{Sport}{Corsa, Nuoto, Ciclismo}
\cvline{Lettura}{Saggistica, Fantascienza}
\section{Volontariato}

\cventry{mar 2019–Presente}
        {Volontario e mentore}
        {Security Without Borders Italia}
        {\url{https://www.securitywithoutborders.it}}
        {}
        {\begin{itemize}
\item Mentoring per giovani professionisti e studenti che si avvicinano alla sicurezza informatica.
\item Organizzazione di laboratori pratici su phishing, igiene digitale e protezione dei dati.
\item Supporto alla community nella gestione responsabile delle vulnerabilità.
\end{itemize}}
    

\cventry{ott 2017–dic 2018}
        {Formatore digitale}
        {Associazione Informatica Solidale}
        {\url{https://www.informaticasolidale.it}}
        {}
        {\begin{itemize}
\item Corsi di alfabetizzazione digitale e sicurezza online per persone over 60.
\item Creazione di materiali divulgativi semplici e accessibili.
\item Collaborazione con biblioteche e centri civici del territorio.
\end{itemize}}
    
\end{document}